\subsection{Mismatch/Outlier-Aware Boundary Compensation}
\label{sec:method_boundary_comp}

AMCPrune first removes a contiguous redundant depth interval and then repairs the newly exposed boundary with a lightweight channel-selective mapping. Let $B_i$ and $B_j$ denote the left and right boundary blocks selected by the depth-pruning stage, where the intermediate blocks $B_{i+1},\ldots,B_{j-1}$ are physically removed. In the dense model, $B_j$ receives hidden states produced by $B_{j-1}$, while in the pruned model it directly receives the output of $B_i$. This replacement creates a boundary distribution shift even when the removed interval has high hidden-state similarity. We therefore measure the shift at the hidden-channel level and compensate only the most affected channels.

Given calibration inputs, we collect the dense boundary input $H^{\mathrm{dense}}\in\mathbb{R}^{T\times C}$ and the pruned boundary input $H^{\mathrm{pruned}}\in\mathbb{R}^{T\times C}$, where $T$ is the number of calibration tokens and $C$ is the hidden dimension. For each channel $c$, we define the boundary mismatch score
\begin{equation}
m_c =
\left|\mu_c^{\mathrm{dense}}-\mu_c^{\mathrm{pruned}}\right|
+ \lambda_{\sigma}
\left|\sigma_c^{\mathrm{dense}}-\sigma_c^{\mathrm{pruned}}\right|,
\end{equation}
where $\mu_c$ and $\sigma_c$ are the token-wise mean and standard deviation of channel $c$. This term identifies channels whose input statistics change after depth pruning.

Depth pruning can also amplify sensitive activation channels. To account for this, we use a lightweight outlier-risk score
\begin{equation}
o_c =
\mathbb{E}_{t}\left[
|z_{t,c}^{\mathrm{pruned}}|
\cdot
\mathbf{1}\left(|z_{t,c}^{\mathrm{pruned}}|>\tau\right)
\right],
\end{equation}
where $z_{t,c}^{\mathrm{pruned}}$ is the standardized activation of channel $c$ at token $t$, and $\tau$ is an outlier threshold. The final channel-selection score is
\begin{equation}
s_c =
\mathrm{Norm}(m_c)
+ \lambda_o \mathrm{Norm}(o_c),
\qquad
\mathcal{S}=\mathrm{TopK}_{c}(s_c,\rho C),
\end{equation}
where $\rho$ is the compensation channel ratio. Unlike normalization-only correction, this selection explicitly decides which hidden channels require boundary repair.

For the selected channels, we first estimate a channel-wise affine mapping from calibration pairs:
\begin{equation}
\hat{H}_{t,c} =
\begin{cases}
\alpha_c H^{\mathrm{pruned}}_{t,c} + \beta_c, & c\in\mathcal{S},\\
H^{\mathrm{pruned}}_{t,c}, & c\notin\mathcal{S}.
\end{cases}
\end{equation}
This affine correction is parameter-efficient, but it cannot model cross-channel interactions introduced by the removed attention and MLP sublayers. We therefore use the affine variant as a diagnostic baseline and introduce a stronger residual boundary compensator:
\begin{equation}
\hat{H}
= H^{\mathrm{pruned}}
+ \gamma U \phi(V H^{\mathrm{pruned}}),
\qquad
V\in\mathbb{R}^{r\times C},\quad
U\in\mathbb{R}^{C\times r},
\end{equation}
where $r\ll C$ is the adapter rank, $\phi$ is a GELU nonlinearity, and $\gamma$ controls the residual strength. The adapter is trained only on the calibration boundary pairs by minimizing the reconstruction loss on the selected channels:
\begin{equation}
\mathcal{L}_{\mathrm{comp}}
=
\frac{1}{T|\mathcal{S}|}
\sum_{t=1}^{T}\sum_{c\in\mathcal{S}}
\left\|
\hat{H}_{t,c}
- H^{\mathrm{dense}}_{t,c}
\right\|_2^2 .
\end{equation}
The compensation module is inserted only at the pruned boundary before the right boundary block $B_j$. This keeps the compact model structurally close to the depth-pruned model while giving the compensation function enough capacity to recover boundary transfer behavior that channel-wise affine rescaling cannot express.

\paragraph{Ablation design.}
To isolate the role of channel selection and compensation capacity, we compare random, mismatch-only, outlier-only, and mismatch+outlier channel selection under the same depth-pruned backbone and channel ratio. We also compare channel-wise affine compensation with the low-rank residual compensator. This ablation tests whether performance gains come from the proposed channel criterion, the residual compensation capacity, or their interaction.
